\textbf{Phần 1. Trắc nghiệm lựa chọn}
\setcounterpageref{ex}{0}

\begin{ex} %[3D3V66]
Những tháng nào trong năm có đúng 31 ngày?
\choice
{Tháng 1, 2, 3, 4, 5, 6, 7}
{Tháng 4, 6, 9, 11}
{Tháng 1, 3, 5, 7, 9, 11}
{\True Tháng 1, 3, 5, 7, 8, 10, 12}
\loigiai{
Theo quy tắc các khớp ngón tay: tháng 1, 3, 5, 7, 8, 10, 12 có đúng 31 ngày.
\textbf{Đáp án của câu hỏi là D}
}
\end{ex}

\begin{ex} %[3D3V66]
Nếu ngày 3 tháng 10 là ngày Chủ nhật thì ngày 24 tháng 10 cùng năm đó là ngày thứ mấy?
\choice
{Thứ Bảy}
{\True Chủ nhật}
{Thứ Hai}
{Thứ Sáu}
\loigiai{
Khoảng cách: $24 - 3 = 21$ ngày. Vì $21 : 7 = 3$ (chia hết cho 7) nên ngày 24 tháng 10 cũng là ngày Chủ nhật.
\textbf{Đáp án của câu hỏi là B}
}
\end{ex}

\begin{ex} %[3D3V66]
Đồng hồ chỉ kim ngắn ở giữa số 4 và số 5, kim dài chỉ số 8. Lúc đó là mấy giờ?
\choice
{$4$ giờ $8$ phút}
{$5$ giờ $40$ phút}
{\True $4$ giờ $40$ phút (hoặc $5$ giờ kém $20$ phút)}
{$4$ giờ $20$ phút}
\loigiai{
Kim dài chỉ số 8 là: $8 \times 5 = 40$ phút. Lúc đó là 4 giờ 40 phút hay 5 giờ kém 20 phút.
\textbf{Đáp án của câu hỏi là C}
}
\end{ex}

\begin{ex} %[3D3V66]
Tháng 2 của năm nhuận có bao nhiêu ngày?
\choice
{\True $29$ ngày}
{$28$ ngày}
{$30$ ngày}
{$31$ ngày}
\loigiai{
Năm nhuận có 366 ngày, riêng tháng 2 có 29 ngày.
\textbf{Đáp án của câu hỏi là A}
}
\end{ex}

\begin{ex} %[3D3V66]
Một bộ phim bắt đầu chiếu lúc $19$ giờ $30$ phút và kết thúc lúc $21$ giờ $15$ phút. Bộ phim đó kéo dài bao lâu?
\choice
{$1$ giờ $30$ phút}
{$2$ giờ}
{$1$ giờ $15$ phút}
{\True $1$ giờ $45$ phút}
\loigiai{
Từ 19 giờ 30 phút đến 20 giờ là 30 phút; từ 20 giờ đến 21 giờ 15 phút là 1 giờ 15 phút. Tổng: $30 + 75 = 105$ phút = 1 giờ 45 phút.
\textbf{Đáp án của câu hỏi là D}
}
\end{ex}

\begin{ex} %[3D3V66]
Nếu ngày 28 tháng 4 là thứ Năm thì ngày 2 tháng 5 cùng năm đó là ngày thứ mấy?
\choice
{Chủ nhật}
{\True Thứ Hai}
{Thứ Ba}
{Thứ Bảy}
\loigiai{
Tháng 4 có 30 ngày: ngày 28 (thứ Năm), 29 (thứ Sáu), 30 (thứ Bảy), ngày 1 tháng 5 (Chủ nhật), ngày 2 tháng 5 là thứ Hai.
\textbf{Đáp án của câu hỏi là B}
}
\end{ex}

\begin{ex} %[3D3V66]
Những tháng nào sau đây có đúng 30 ngày?
\choice
{Tháng 2, tháng 4, tháng 6, tháng 8}
{Tháng 1, tháng 4, tháng 7, tháng 10}
{\True Tháng 4, tháng 6, tháng 9, tháng 11}
{Tháng 3, tháng 5, tháng 9, tháng 11}
\loigiai{
Bốn tháng có đúng 30 ngày là: tháng 4, 6, 9, 11.
\textbf{Đáp án của câu hỏi là C}
}
\end{ex}

\begin{ex} %[3D3V66]
Một đồng hồ chạy chậm mỗi ngày 3 phút. Hỏi sau một tuần lễ, chiếc đồng hồ đó chạy chậm tất cả bao nhiêu phút?
\choice
{\True $21$ phút}
{$15$ phút}
{$18$ phút}
{$24$ phút}
\loigiai{
Một tuần lễ có 7 ngày: $3 \times 7 = 21$ phút.
\textbf{Đáp án của câu hỏi là A}
}
\end{ex}

\begin{ex} %[3D3V66]
Nam bắt đầu học bài lúc $20$ giờ $10$ phút và học xong lúc $21$ giờ $05$ phút. Thời gian Nam ngồi học là
\choice
{$45$ phút}
{$50$ phút}
{$1$ giờ}
{\True $55$ phút}
\loigiai{
Từ 20 giờ 10 phút đến 21 giờ là 50 phút; thêm 5 phút nữa là: $50 + 5 = 55$ phút.
\textbf{Đáp án của câu hỏi là D}
}
\end{ex}

\begin{ex} %[3D3V66]
Năm không nhuận có tất cả bao nhiêu ngày?
\choice
{$366$ ngày}
{\True $365$ ngày}
{$360$ ngày}
{$364$ ngày}
\loigiai{
Năm thường (không nhuận) có 365 ngày.
\textbf{Đáp án của câu hỏi là B}
}
\end{ex}

\begin{ex} %[3D3V66]
Lúc $15$ giờ $45$ phút còn được gọi là
\choice
{$3$ giờ $45$ phút sáng}
{$4$ giờ $15$ phút chiều}
{\True $4$ giờ kém $15$ phút chiều}
{$15$ giờ kém $15$ phút}
\loigiai{
$15$ giờ là 3 giờ chiều. $15$ giờ $45$ phút chính là 4 giờ kém 15 phút chiều.
\textbf{Đáp án của câu hỏi là C}
}
\end{ex}

\begin{ex} %[3D3V66]
Một trận bóng đá diễn ra trong hai hiệp, mỗi hiệp $45$ phút, nghỉ giữa hai hiệp $15$ phút. Tổng thời gian của trận bóng đá (không kể bù giờ) là
\choice
{\True $1$ giờ $45$ phút}
{$1$ giờ $30$ phút}
{$2$ giờ}
{$1$ giờ $15$ phút}
\loigiai{
Tổng thời gian: $45 \times 2 + 15 = 90 + 15 = 105$ phút = 1 giờ 45 phút.
\textbf{Đáp án của câu hỏi là A}
}
\end{ex}

\textbf{Phần 2. Trắc nghiệm đúng sai}
\setcounterpageref{ex}{0}

\begin{ex} %[3D3V66]
Xét tính đúng sai của các nhận định về thời gian, tháng và năm:
\choiceTF
{\True Một năm có 12 tháng.}
{Tất cả các tháng có 31 ngày đều là các tháng mang số lẻ.}
{\True Tháng 2 năm nhuận có 29 ngày, năm thường có 28 ngày.}
{Một ngày có 12 giờ.}
\loigiai{
\begin{itemize}
    \item Ý a) Đúng: Kiến thức chuẩn.
    \item Ý b) Sai: Sai vì tháng 8, 10, 12 là số chẵn nhưng vẫn có 31 ngày.
    \item Ý c) Đúng: Quy luật lịch dương.
    \item Ý d) Sai: Sai vì một ngày có 24 giờ.
\end{itemize}
\begin{itemize}
    \item a) \textbf{Đúng}.
    \item b) \textbf{Sai}.
    \item c) \textbf{Đúng}.
    \item d) \textbf{Sai}.
\end{itemize}
}
\end{ex}

\begin{ex} %[3D3V66]
Lịch tháng 8 của một năm có ngày 19 tháng 8 là ngày thứ Sáu:
\choiceTF
{\True Ngày 12 tháng 8 cùng năm đó cũng là ngày thứ Sáu.}
{\True Ngày 26 tháng 8 cùng năm đó cũng là ngày thứ Sáu.}
{Tháng 8 có đúng 30 ngày.}
{\True Ngày 2 tháng 9 cùng năm đó là ngày thứ Sáu.}
\loigiai{
\begin{itemize}
    \item Ý a) Đúng: $19 - 7 = 12$.
    \item Ý b) Đúng: $19 + 7 = 26$.
    \item Ý c) Sai: Sai vì tháng 8 có 31 ngày.
    \item Ý d) Đúng: Tháng 8 có 31 ngày: 26 (thứ Sáu) + 7 ngày = ngày 2 tháng 9 (thứ Sáu).
\end{itemize}
\begin{itemize}
    \item a) \textbf{Đúng}.
    \item b) \textbf{Đúng}.
    \item c) \textbf{Sai}.
    \item d) \textbf{Đúng}.
\end{itemize}
}
\end{ex}

\textbf{Phần 3. Trả lời ngắn}
\setcounterpageref{ex}{0}

\begin{ex} %[3D3V66]
Một năm thường (không nhuận) có bao nhiêu ngày?
\shortans{365}
\loigiai{
\textbf{Đáp số:} $365$\\
Năm thường có 365 ngày.
}
\end{ex}

\begin{ex} %[3D3V66]
Có bao nhiêu tháng trong năm có đúng 30 ngày?
\shortans{4}
\loigiai{
\textbf{Đáp số:} $4$\\
Gồm 4 tháng: 4, 6, 9, 11.
}
\end{ex}

\begin{ex} %[3D3V66]
Nếu ngày 5 tháng 5 là thứ Năm thì ngày 26 tháng 5 cùng năm là thứ mấy (viết số thứ: thứ 5 viết 5)?
\shortans{5}
\loigiai{
\textbf{Đáp số:} $5$\\
$26 - 5 = 21$ chia hết cho 7 nên vẫn là thứ Năm (5).
}
\end{ex}

\begin{ex} %[3D3V66]
Đổi $2$ giờ $15$ phút ra đơn vị phút.
\shortans{135}
\loigiai{
\textbf{Đáp số:} $135$\\
$2 \times 60 + 15 = 135$ phút.
}
\end{ex}

\begin{ex} %[3D3V66]
Xe ô tô khởi hành lúc 7 giờ 30 phút và đến nơi lúc 10 giờ 15 phút. Hỏi xe chạy hết bao nhiêu phút?
\shortans{165}
\loigiai{
\textbf{Đáp số:} $165$\\
Từ 7h30 đến 10h15 là 2 giờ 45 phút = $120 + 45 = 165$ phút.
}
\end{ex}

\begin{ex} %[3D3V66]
Tháng 12 có bao nhiêu ngày?
\shortans{31}
\loigiai{
\textbf{Đáp số:} $31$\\
Tháng 12 có 31 ngày.
}
\end{ex}
